% Options for packages loaded elsewhere
\PassOptionsToPackage{unicode}{hyperref}
\PassOptionsToPackage{hyphens}{url}
\documentclass[
]{article}
\usepackage{xcolor}
\usepackage{amsmath,amssymb}
\setcounter{secnumdepth}{-\maxdimen} % remove section numbering
\usepackage{iftex}
\ifPDFTeX
  \usepackage[T1]{fontenc}
  \usepackage[utf8]{inputenc}
  \usepackage{textcomp} % provide euro and other symbols
\else % if luatex or xetex
  \usepackage{unicode-math} % this also loads fontspec
  \defaultfontfeatures{Scale=MatchLowercase}
  \defaultfontfeatures[\rmfamily]{Ligatures=TeX,Scale=1}
\fi
\usepackage{lmodern}
\ifPDFTeX\else
  % xetex/luatex font selection
\fi
% Use upquote if available, for straight quotes in verbatim environments
\IfFileExists{upquote.sty}{\usepackage{upquote}}{}
\IfFileExists{microtype.sty}{% use microtype if available
  \usepackage[]{microtype}
  \UseMicrotypeSet[protrusion]{basicmath} % disable protrusion for tt fonts
}{}
\makeatletter
\@ifundefined{KOMAClassName}{% if non-KOMA class
  \IfFileExists{parskip.sty}{%
    \usepackage{parskip}
  }{% else
    \setlength{\parindent}{0pt}
    \setlength{\parskip}{6pt plus 2pt minus 1pt}}
}{% if KOMA class
  \KOMAoptions{parskip=half}}
\makeatother
\usepackage{color}
\usepackage{fancyvrb}
\newcommand{\VerbBar}{|}
\newcommand{\VERB}{\Verb[commandchars=\\\{\}]}
\DefineVerbatimEnvironment{Highlighting}{Verbatim}{commandchars=\\\{\}}
% Add ',fontsize=\small' for more characters per line
\newenvironment{Shaded}{}{}
\newcommand{\AlertTok}[1]{\textcolor[rgb]{1.00,0.00,0.00}{\textbf{#1}}}
\newcommand{\AnnotationTok}[1]{\textcolor[rgb]{0.38,0.63,0.69}{\textbf{\textit{#1}}}}
\newcommand{\AttributeTok}[1]{\textcolor[rgb]{0.49,0.56,0.16}{#1}}
\newcommand{\BaseNTok}[1]{\textcolor[rgb]{0.25,0.63,0.44}{#1}}
\newcommand{\BuiltInTok}[1]{\textcolor[rgb]{0.00,0.50,0.00}{#1}}
\newcommand{\CharTok}[1]{\textcolor[rgb]{0.25,0.44,0.63}{#1}}
\newcommand{\CommentTok}[1]{\textcolor[rgb]{0.38,0.63,0.69}{\textit{#1}}}
\newcommand{\CommentVarTok}[1]{\textcolor[rgb]{0.38,0.63,0.69}{\textbf{\textit{#1}}}}
\newcommand{\ConstantTok}[1]{\textcolor[rgb]{0.53,0.00,0.00}{#1}}
\newcommand{\ControlFlowTok}[1]{\textcolor[rgb]{0.00,0.44,0.13}{\textbf{#1}}}
\newcommand{\DataTypeTok}[1]{\textcolor[rgb]{0.56,0.13,0.00}{#1}}
\newcommand{\DecValTok}[1]{\textcolor[rgb]{0.25,0.63,0.44}{#1}}
\newcommand{\DocumentationTok}[1]{\textcolor[rgb]{0.73,0.13,0.13}{\textit{#1}}}
\newcommand{\ErrorTok}[1]{\textcolor[rgb]{1.00,0.00,0.00}{\textbf{#1}}}
\newcommand{\ExtensionTok}[1]{#1}
\newcommand{\FloatTok}[1]{\textcolor[rgb]{0.25,0.63,0.44}{#1}}
\newcommand{\FunctionTok}[1]{\textcolor[rgb]{0.02,0.16,0.49}{#1}}
\newcommand{\ImportTok}[1]{\textcolor[rgb]{0.00,0.50,0.00}{\textbf{#1}}}
\newcommand{\InformationTok}[1]{\textcolor[rgb]{0.38,0.63,0.69}{\textbf{\textit{#1}}}}
\newcommand{\KeywordTok}[1]{\textcolor[rgb]{0.00,0.44,0.13}{\textbf{#1}}}
\newcommand{\NormalTok}[1]{#1}
\newcommand{\OperatorTok}[1]{\textcolor[rgb]{0.40,0.40,0.40}{#1}}
\newcommand{\OtherTok}[1]{\textcolor[rgb]{0.00,0.44,0.13}{#1}}
\newcommand{\PreprocessorTok}[1]{\textcolor[rgb]{0.74,0.48,0.00}{#1}}
\newcommand{\RegionMarkerTok}[1]{#1}
\newcommand{\SpecialCharTok}[1]{\textcolor[rgb]{0.25,0.44,0.63}{#1}}
\newcommand{\SpecialStringTok}[1]{\textcolor[rgb]{0.73,0.40,0.53}{#1}}
\newcommand{\StringTok}[1]{\textcolor[rgb]{0.25,0.44,0.63}{#1}}
\newcommand{\VariableTok}[1]{\textcolor[rgb]{0.10,0.09,0.49}{#1}}
\newcommand{\VerbatimStringTok}[1]{\textcolor[rgb]{0.25,0.44,0.63}{#1}}
\newcommand{\WarningTok}[1]{\textcolor[rgb]{0.38,0.63,0.69}{\textbf{\textit{#1}}}}
\setlength{\emergencystretch}{3em} % prevent overfull lines
\providecommand{\tightlist}{%
  \setlength{\itemsep}{0pt}\setlength{\parskip}{0pt}}
\usepackage{bookmark}
\IfFileExists{xurl.sty}{\usepackage{xurl}}{} % add URL line breaks if available
\urlstyle{same}
\hypersetup{
  hidelinks,
  pdfcreator={LaTeX via pandoc}}

\author{}
\date{}

\begin{document}

\section{Methods Included: Standardizing Computational Reuse and
Portability with the Common Workflow
Language}\label{methods-included-standardizing-computational-reuse-and-portability-with-the-common-workflow-language}

\section{Michael R. Crusoe}\label{michael-r.-crusoe}

VU Amsterdam

Department of Computer Science

Amsterdam, The Netherlands

Software Freedom Conservancy

Common Workflow Language project

Brooklyn, NY, USA

mrc@commonwl.org

\section{Peter Amstutz}\label{peter-amstutz}

Curii Corporation

Sommerville, MA, USA

peter.amstutz@curii.com

\section{Sanne Abeln}\label{sanne-abeln}

VU Amsterdam

Department of Computer Science

Amsterdam, The Netherlands

s.abeln@vu.nl

\section{Alexandru Iosup}\label{alexandru-iosup}

VU Amsterdam

Department of Computer Science

Amsterdam, The Netherlands

a.iosup@vu.nl

\section{Hervé Ménager}\label{hervuxe9-muxe9nager}

Institut Pasteur

Hub de Bioinformatique et

Biostatistique -- Département Biologie

Computationnelle

Paris, France

herve.menager@pasteur.fr

\section{John Chilton}\label{john-chilton}

Pennsylvania State University

Department of Biochemistry and

Molecular Biology

State College, PA, USA

Galaxy Project

Multiple Countries

jmchilton@gmail.com

\section{Stian Soiland-Reyes}\label{stian-soiland-reyes}

The University of Manchester

Department of Computer Science

Manchester, UK

Informatics Institute, University of

Amsterdam

Amsterdam, Netherlands

soiland-reyes@manchester.ac.uk

\section{Nebojša Tijanić}\label{nebojux161a-tijaniux107}

Seven Bridges

Charlestown, MA, USA

boysha@gmail.com

\section{Bogdan Gavrilovic}\label{bogdan-gavrilovic}

Seven Bridges

Charlestown, MA, USA

bogdan.gavrilovic@sbgenomics.com

\section{Carole Goble}\label{carole-goble}

The University of Manchester

Department of Computer Science

Manchester, UK

carole.goble@manchester.ac.uk

\section{The CWL Community}\label{the-cwl-community}

Software Freedom Conservancy

Common Workflow Language project

Brooklyn, NY, USA

common-workflow-

language@googlegroups.com

\section{ABSTRACT}\label{abstract}

Computational Workflows are widely used in data analysis, enabling
innovation and decision-making for the modern society. In many domains
the analysis components are numerous and written in multiple different
computer languages by third parties. These polylingual workflows are
common in many industries and dominant in fields such as bioinformatics,
image analysis, and radio astronomy. However, in practice many competing
workflow systems exist, severely limiting portability of such workflows,
thereby hindering the transfer of such workflows between different
systems, between different projects and different settings, leading to
vendor lock-ins and limiting their generic re-usability. Here we present
the Common Workflow Language (CWL) project which produces free and open
standards for describing command-line tool based workflows. The CWL
standards provide a common but

reduced set of abstractions that are both used in practice and
implemented in many popular workflow systems. The CWL language is
declarative, which allows expressing computational workflows constructed
from diverse software tools, executed each through their command-line
interface. Being explicit about the runtime environment and any use of
software containers enables portability and reuse. The CWL project is
not specific to a particular analysis domain, it is community-driven,
and it produces consensus-built standards. Workflows written according
to the CWL standards are a reusable description of that analysis that
are runnable on a diverse set of computing environments. These
descriptions contain enough information for advanced optimization
without additional input from workflow authors. The CWL standards
support polylingual workflows, enabling portability and reuse of such
workflows, easing for example scholarly publication, fulfilling
regulatory requirements, collaboration in/between academic research and
industry,

while reducing implementation costs. CWL has been taken up by a wide
variety of domains, and industries and support has been implemented in
many major workflow systems.

\section{CCS CONCEPTS}\label{ccs-concepts}

- Computing methodologies → Distributed computing methodologies; •
General and reference → Computing standards, RFCs and guidelines; •
Applied computing → Astronomy; Earth and atmospheric sciences;
Enterprise interoperability; Enterprise computing infrastructures; Life
and medical sciences; Bioinformatics; Transcriptomics; Computational
proteomics; Population genetics; Systems biology; Computational biology;
Computational proteomics; Computational genomics; Imaging; Computational
transcriptomics; • Computer systems organization → Distributed
architectures; Cloud computing; Grid computing; Data flow architectures.

\section{KEYWORDS}\label{keywords}

workflows, computational data analysis, CWL, scientific workflows,
standards

\section{1 INTRODUCTION}\label{introduction}

Computational Workflows are widely used in data analysis, enabling
innovation and decision-making for the modern society. But their growing
popularity is also a cause for concern: unless we standardize
computational reuse and portability, the use of workflows may end up
hampering collaboration. How can we enjoy the common benefits of
computational workflows and also eliminate such risks?

To answer this general question, we advocate in this work for workflow
thinking as a shared way of reasoning across all domains and
practitioners, introduce Common Workflow Language (CWL) as a pragmatic
set of standards for describing and sharing computational workflows, and
discuss the principles around which these standards have become central
to a diverse community of users across multiple fields of science and
engineering. This article focuses on an overview of the CWL standards
and the CWL project and is complemented by the technical detail
available in the CWL standards themselves \(^{1}\) .

Workflow thinking is a form of ``conceptualizing processes as recipes
and protocols, structured as {[}work- or{]} dataflow graphs with
computational steps, and subsequently developing tools and approaches
for formalizing, analyzing and communicating these process
descriptions'' {[}15{]}. It introduces an abstraction, the workflow,
which helps decouple expertise in a specific domain, for example of
specific science or engineering fields, from expertise in computing.
Derived from workflow thinking, a computational workflow describes a
process for computing where different parts of the process (the tasks)
are inter-dependent, e.g., a task can start processing after its
predecessors have (partially) completed and where data flows between
tasks.

Currently, many competing systems exist that enable simple workflow
execution (workflow runners) or offer a comprehensive management of
workflows and data (workflow management systems), each with their own
syntax or method for describing workflows and infrastructure
requirements. This limits computational reuse and portability. In
particular, although the data-flows are becoming increasingly more
complex, most workflow abstractions do not enable explicit
specifications of data-flows, increasing significantly the costs to
reuse and port the workflow by third-parties.

We thus identify an important problem for the broad adoption of workflow
thinking in practice: although communities require polylingual workflows
(workflows that execute tools written in multiple different computer
languages) and multi-party workflows, adopting and managing different
workflow systems is costly and difficult. In this work, we propose to
tame this complexity through a common abstraction that covers the
majority of features used in practice, and is (or can be) implemented in
many workflow systems.

In the computational workflow depicted in Figure 1, practitioners solved
the problem by adopting the CWL standards. We posit in this work that
the CWL standards provide the common abstraction that can help solve the
main problems of sharing workflows between institutions and users. CWL
achieves this by providing a declarative language that allows expressing
computational workflows constructed from diverse software tools,
executed each through their command-line interface, with the inputs and
outputs of each tool clearly specified and with inputs possibly
resulting from the execution of other tools. We also set out to
introduce the CWL standards, with a tri-fold focus: (1) the CWL
standards focuses on maintaining a separation of concerns between the
description and execution of tools and workflows, proposing a language
that includes only operations commonly used across multiple communities
of practice; (2) the CWL standards support workflow automation,
scalability, abstraction, provenance, portability, and reusability; and
(3) the CWL project takes a principled, community-first open-source and
open-standard approach which enables this result.

The CWL standards are the product of an open and free standards-making
community. While the CWL project began in the bioinformatics domain the
many contributors to the CWL project shaped the standards so that it
could be useful anywhere that experiences the problem of ``many tools
written in many programming languages by many parties''. Since the
ratification of the first version in 2016, the CWL standards have been
used in other fields including hydrology \(^{2}\) , radio astronomy
\(^{3}\) , geo-spatial analysis \([13, 22, 30]\) , high energy physics
\([5]\) , in addition to fast-growing bioinformatics fields like
(meta-)genomics \([26]\) and cancer research \([23]\) . The CWL
standards are featured in US FDA sponsored and adopted IEEE Std 2791
\(^{™}\) -2020 standard \([1]\) and the Netherlands' National Plan for
Open Science \([32]\) . A list of free and open-source implementations
of the CWL standards are listed in Table 1. Additionally there are
multiple commercially supported systems that support the CWL standards
for executing workflows and they are available from vendors such as
Curii (Arvados) \(^{4}\) , DNAnexus \(^{5}\) , IBM (IBM \(^{®}\)
Spectrum LSF) \(^{6}\) , Illumina (Illumina Connected Analytics)
\(^{7}\) , and Seven

\textit{[Image omitted: images/a9deda6898a336cdd461c824048c46c7511ccd54b7027df1a428b8d527489488.jpg]}

flowchart

Workflow inputs and outputs flowchart showing data flow between
individual workflows, sequence sequences, and specific matches like 16S,
5S, 23S, and tRNA.

Figure 1: Excerpt from a large microbiome bioinformatics CWL workflow
{[}26{]}. This part of the workflow (which is interpretable/executable
on its own) has the aim to match the workflow inputs of genomic
sequences to provided sequence-models, which are dispatched to four
sub-workflows (e.g., find\_16S\_matches); the sub-workflows not detailed
in the figure. The sub-worklow outputs are then collated to identify
unique sequence hits, then provided as overall workflow outputs. Arrows
define the connection between tasks and imply their partial ordering,
depicted here as layers of tasks that may execute concurrently. Workflow
steps (e.g., mask\_rRNA\_and\_tRNA) execute command line tools, shown
here with indicators for their different programming languages (e.g.,
{[}Py{]} for Python, {[}C{]} for the C language). (Diagram adapted from
https://w3id.org/cwl/view/git/7bb76f33bf40b5cd2604001cac46f967a209c47f/workflows/rna-selector.cwl,
which was originally retrieved from a corresponding CWL workflow of the
EBI Metagenomics project, itself a conversion of the ``rRNASelector''
{[}24{]} program into a well structured workflow allowing for better
parallelization of execution and provenance tracking.)

Bridges \(^{8}\) . The flexibility of the CWL standards enabled, for
example, rapid collaboration on and prototyping of a COVID-19 public
database and analysis resource {[}16{]}.

The separation of concerns proposed by the CWL standards enable diverse
projects, and can also benefit engineering and large industrial
projects. Likewise, users of Docker (or other software container
technologies) that distribute analysis tools can use just the CWL
Command Line Tool standard for providing a structured
workflow-independent description of how to run their tool(s) in the
container, what data is required to be provided to the container, and
what results to expect and where to find them in the container.

\section{Key Insights {[}in CACM box{]}}\label{key-insights-in-cacm-box}

Toward computational reuse and portability of polylingual, multi-party
workflows, the CWL project makes the following contributions:

\begin{enumerate}
\def\labelenumi{(\arabic{enumi})}
\item
  CWL is a set of standards for describing and sharing computational
  workflows.\\
\item
  The CWL standards are used daily in many science and engineering
  domains, including by multi-stakeholder teams.\\
\item
  The CWL standards use a declarative syntax, facilitating polylingual
  workflow tasks. By being explicit about the runtime environment and
  any use of software containers, the CWL standards enable portability
  and reuse. (See Section 3.)\\
\item
  The CWL standards provide a separation of concerns between workflow
  authors and workflow platforms. (More in Section 4.3.)
\item
  The CWL standards support critical workflow concepts like automation,
  scalability, abstraction, provenance, portability, and reusability.
  (Details in Section 2).\\
\item
  The CWL standards are developed around core principles of community
  and shared decision-making, re-use, and zero cost for participants.
  (Section 4 details the open standards.)\\
\item
  The CWL standards are provided as freely available open standards,
  supported by a diverse community in collaboration with industry, and
  is a Free/Open Source Software ecosystem (see Sidebar B, Section 4.2).
\end{enumerate}

\section{2 BACKGROUND ON WORKFLOWS AND STANDARDS FOR
WORKFLOWS}\label{background-on-workflows-and-standards-for-workflows}

Workflows, and standards-based descriptions thereof, hold the potential
to solve key problems in many domains of science and engineering. This
section explains why.

\section{2.1 Why Workflows?}\label{why-workflows}

In many domains, workflows include diverse analysis components, written
in multiple (different) computer languages, by both end-users and
third-parties. Such polylingual and multi-party workflows are already
common or dominant in data-intensive fields like bioinformatics, image
analysis, and radio astronomy; we envision they could bring important
benefits to many other domains.

To thread data through analysis tools, domain experts such as
bioinformaticians use specialized command-line interfaces {[}12, 29{]}
and other domains use their own customized frameworks {[}3, 6{]}.

Workflow engines also help with efficient management of the resources
used to run scientific workloads {[}7, 10{]}.

The workflow approach helps compose an entire application of these
command-line analysis tools: developers build graphical or textual
descriptions of how to run these command-line tools, and scientists and
engineers connect their inputs and outputs so that the data flows
through. An example of a complex workflow problem is metagenomic
analysis, for which Figure 1 illustrates a subset (a sub-workflow).

In practice, many research and engineering groups use workflows of the
kind described in Figure 1. However, as highlighted in a recently
published ``Technology Toolbox'' article {[}28{]} published in the
journal Nature, these groups typically lack the ability to share and
collaborate across institutions and infrastructures without costly
manual translation of their workflows.

Using workflow techniques, especially with digital analysis processes,
has become quite popular and does not look to be slowing down: one
workflow management system recently celebrated its 10,000th citation
\(^{9}\) ; and over 298 computational data analysis workflow systems are
known \(^{10}\) .

A process, digital or otherwise, may grow to such complexity that the
authors and users of that process have difficulties in understanding its
structure, scaling the process, managing the running of the process, and
keeping track of what happened in previous enactments of the process.
Process dependencies may be undocumented, obfuscated, or otherwise
effectively invisible; even an extensively documented process may be
difficult to understand by outsiders or newcomers if a common framework
or vocabulary is lacking. The need to run the process more frequently or
with larger inputs is unlikely to be achieved by the initial entity
(i.e., either script or person) running the process. What seemed once a
reasonable manual step (run this command here and then paste the result
there; then call this person for permission) will, under the pressure of
porting and reusing, become a bottleneck. Informal logs (if any) will
quickly become unsuitable for answering an organization's need to
understand what happened, when, by whom, and to which data.

Workflow techniques aim to solve these problems by providing the
Abstraction, Scaling, Automation, and Provenance (A.S.A.P.) features
{[}8{]}. Workflow constructs enable a clear abstraction about the
components, the relationships between components, and the inputs and
outputs of the components turning them into well-labeled tools with
documented expectations. This abstraction enables scaling (execution can
be parallelized and distributed), automation (the abstraction can be
used by a workflow engine to track, plan, and manage execution of
tasks), and provenance tracking (descriptions of tasks, executors,
inputs, outputs; with timestamps, identifiers (unique names), and other
logs, can be stored in relation to each other to later answer structured
queries).

\section{2.2 Why Workflow Standards?}\label{why-workflow-standards}

Although workflows are very popular, prior to the CWL standards every
workflow system was incompatible with every other. This means that those
users not using the CWL standards are required to express their
computational workflows in a different way every time they have to use
another workflow system leading to local success, but global
unportability.

The success of workflows is now their biggest drawback: users are locked
into a particular vendor, project, and often a particular hardware
setup. This hampers sharing and re-use. Even non-academics suffer from
this situation, as the lack of standards (or the lack of their adoption)
hinders effective collaboration on computational methods within and
between companies. Likewise, this unportability affects public-private
partnerships and the potential for technology transfer from public
researchers.

A second significant problem is that incomplete method descriptions are
common when computational analysis is reported in academic research
{[}17{]}. Reproduction, re-use, and replication {[}11{]} of these
digital methods requires a complete description of what computer
applications were used, how exactly they were used, and how they were
connected to each other. For precision and interoperability, this
description should also be in an appropriate standardized
machine-readable format.

A standard for sharing and reusing workflows can provide a solution to
describing portable, re-usable workflows while also being
workflow-engine and vendor-neutral.

Sharing workflow descriptions based on standards also addresses the
second problem: the availability of the workflow description provides
needed information when sharing; and the quality of the description
provided by a structured, standards-based approach is much higher than
the current approach of casual, unstructured, and almost always
incomplete descriptions in scientific reports. Moreover, the operational
parts of the description can be provided automated by the workflow
management system, rather than by domain experts.

While (data) standards are commonly adopted and have become expected for
funded projects in knowledge representation fields, the same cannot yet
be said about workflows and workflow engines yet.

\section{2.3 Sidebar A: Monolingual and Polylingual workflow
systems}\label{sidebar-a-monolingual-and-polylingual-workflow-systems}

Workflows techniques can be implemented in many ways, i.e., with varying
degrees of formalism, which tends to correlate with execution
flexibility and features. Typically, whereas the most informal
techniques require that all processing components are written in the
same programming language or are at least callable from the same
programming language, the formal workflow techniques tend to allow
components to be developed in multiple programming languages.

Among the informal techniques, the do-it-yourself approach uses from a
particular programming language its built-in capabilities. For example,
Python provides a threading library, and the Java-based Apache Hadoop
{[}31{]} provides MapReduce capabilities. To gain more flexibility when
working with a particular programming language, general third-party
libraries, such as ipyparallel \(^{11}\) , can enable remote or
distributed execution without having to re-write one's code.

A more explicit workflow structure can be achieved by using a workflow
library focusing on a specific programming language.

For example, in Parsl {[}3{]}, the workflow constructs (``this is a unit
of processing'', ``here are the dependencies between the units'') are
made explicit and added by the developer to a Python script, to upgrade
it to a scalable workflow. (While we list Parsl here as an example of a
monolingual workflow system, it also contains explicit support for
executing external command-line tools.)

Two approaches can accommodate polylingual workflows where the
components are written in more than one programming language, or where
the components come from third-parties and the user does not want to or
cannot modify them: use of per-language add-in libraries or the use of
the Portable Operating System Interface command-line interface (POSIX
CLI) {[}14{]}. The use of per-language add-in libraries entails either
explicit function calls (e.g., using ctypes in Python to call a C
library \(^{12}\) ) or the addition of annotations to the user's
functions, and requires mapping/restricting to a common cross-language
data model.

Essentially all programming languages support the creation of POSIX CLIs
are familiar to many Linux and macOS users; scripts or binaries which
can be invoked on the shell with a set of arguments, reading and writing
files, and executed in a separate process. Choosing the POSIX
command-line interface as the point of coordination means the connection
between components is done by an array of string arguments representing
program options (including paths to data files) along with a
string-based environment variables (key-value pairs). Using the
command-line as a coordination interface has the advantage of not
needing additional implementation in every programming language, but has
the disadvantages of process start-up time and a very simple data model.
(As a polylingual workflow standard, CWL uses the POSIX CLI data model.)

\section{3 FEATURES OF THE COMMON WORKFLOW LANGUAGE
STANDARDS}\label{features-of-the-common-workflow-language-standards}

The Common Workflow Language standards aim to cover the common needs of
users and the commonly implemented features of workflow runners or
platforms. The remainder of this section presents an overview of the CWL
features, how they translate to executing workflows in CWL format, and
where the CWL standards are not helpful.

The CWL standard support polylingual and multi-party workflows, for
which they enable computational reuse and portability (see also the CACM
Box for main features). To do so, each release of the CWL standards has
two \(^{13}\) main components: (1) a standard for describing command
line tools; and (2) a standard for describing workflows that compose
such tool descriptions. The goal of the CWL Command Line Tool
Description Standard \(^{14}\) is to describe how a particular command
line tool works: what are the inputs and parameters and their types; how
to add the correct flags and switches to the command line invocation;
and where to find the output files.

The CWL standards define an explicit language, both in syntax, and in
its data and execution model. Its textual syntax is derived from YAML
\(^{15}\) . This syntax does not restrict the amount of detail; for
example, Figure 2A depicts a simple example with sparse detail, and
Figure 2B depicts the same example but with the execution augmented with
further details. Each input to a tool has a name and a type (e.g., File,
see label 1 in the figure). Authors of tool descriptions are encouraged
to include documentation and labels for all components (i.e., as in
Figure 2B), to enable the automatic generation of helpful visual
depictions and even Graphical User Interfaces for any given CWL
description. Metadata about the tool description authors themselves
encourages attribution of their efforts. As shown in Figure 2B, item 3,
these tool descriptions can contain well-defined hints or mandatory
requirements such as which software container to use or how much compute
resources are required (memory, number of CPU cores, disk space, and/or
the maximum time or deadline to complete the step or entire workflow.)

The CWL execution model is explicit: Each tool's runtime environment is
explicit and any required elements must be specified by the CWL
tool-description author (in contrast to hints, which are optional)
\(^{16}\) . Each tool invocation uses a separate working directory,
populated according to the CWL tool description, e.g., with the input
files explicitly specified by the workflow author. Some applications
expect particular filenames, directory layouts, and environment
variables, and there are additional constructs in the CWL Command Line
Tool standard to satisfy their needs.

The explicit runtime model enables portability, by being explicit about
data locations. As Figure 3 indicates, this enables execution of CWL
workflows on diverse environments as provided by various implementations
of the CWL standards: the local environment of the author-scientist
(e.g., a single desktop computer, laptop, or workstation), a remote
batch production-environment (e.g., a cluster, an entire datacenter, or
even a global multi-datacenter infrastructure), and an on-demand cloud
environment.

The CWL standards explicitly support the use of software container
technologies, such as Docker and Singularity, to enable portability of
the underlying analysis tools. Figure 2B, item 2, illustrates the
process of pulling a Docker container-image from the Quay.io registry;
then, the workflow engine automates the mounting of files and folders
within the container. The container included in the figure has been
developed by a trusted author and is commonly used in the bioinformatics
field with an expectation its results are reproducible. Indeed, the use
of containers can be seen as a confirmation that a tool's execution is
reproducible, when using only its explicitly declared
runtime-environment. Similarly, when distributed execution is desired,
no changes to the CWL tool-description are needed: because the file or
directory inputs are already explicitly defined in the CWL description,
the (distributed) workflow runner can handle (without additional
configuration) both job placement and data routing between compute
nodes.

Via these two features (special handling of data paths; the optional but
recommended use of software containers), the CWL standards enables
portability (execution ``without change''). Although

\textit{[Image omitted: images/26cea472447c53ff420cecc1d1b8c1a0a7d490d678fe9d76bd723ca7b475eba2.jpg]}\\
Figure 2: Example of CWL syntax and progressive enhancement. (A) and (B)
describe the same tool, but (B) is enhanced with additional features:
human-readable documentation; file format identifiers for better
validation of workflow connections; recommended software container image
for more reproducible results and easier software installation;
dynamically specified resource requirements to optimize task scheduling
and resource usage without manual intervention. The resource
requirements are expressed as hints. (C) shows an example of CWL
Workflow syntax, where the underlying tool descriptions (``grep.cwl''
and ``wc.cwl'') are in external files for ease of reuse.

\textit{[Image omitted: images/1ed633177134d733b834fee0573016083fec3406f924b136888aa00364c20975.jpg]}

flowchart

\begin{Shaded}
\begin{Highlighting}[]
\NormalTok{graph TD}
\NormalTok{    A["Authors of CWL tool and workflow descriptions"] {-}{-}\textgreater{} B["CWL"]}
\NormalTok{    A {-}{-}\textgreater{} C["BATCH"]}
\NormalTok{    A {-}{-}\textgreater{} D["CLOUD"]}
\NormalTok{    B {-}{-}\textgreater{} E["Local execution on Linux, macOS, and MS Windows via the CWL reference implementation (cwltool) and Docker/uDocker/Singularity/podman/..."]}
\NormalTok{    C {-}{-}\textgreater{} F["Backends supported by various F/OSS CWL implementations"]}
\NormalTok{    D {-}{-}\textgreater{} F}
\end{Highlighting}
\end{Shaded}

Figure 3: Example of CWL portability. The same workflow description runs
on the scientist's own laptop or single machine, on any batch
production-environment, and on any common public or private cloud. The
CWL standards enable execution-portability by being explicit about data
locations and execution models.

various factors not controllable by software container technology can
affect portability (e.g., variation in the underlying operating system
kernel; variation in processor results), in practice the exact same
software container and data inputs lead to portability without further
adjustment from the user.

To support features that are not in the CWL standards, the CWL standards
define extension points that permit (namespaced) vendor-specific
features in explicitly defined ways. If these extensions do not
fundamentally change how the tool should operate, then they are added to
the hints list and other CWL compatible engines can ignore them.
However, if the extension is required to properly run the tool being
described, e.g., due to the need for some specialized hardware, then the
extension is listed under requirements and CWL compatible engines can
recognize and explicitly declare their inability to execute that CWL
description.

The CWL Workflow Description Standard \(^{17}\) builds upon the CWL
Command Line Tool Standard: it has the same YAML- or JSON-style syntax,
with explicit workflow level inputs, outputs,

and documentation (see Figure 2C). The workflow descriptions consists of
a list of steps, comprised of CWL CommandLineTools or CWL sub-workflows,
each re-exposing their tool's required inputs. Inputs for each step are
connected by referencing the name of either the common workflow inputs
or particular outputs of other steps. The workflow outputs expose
selected outputs from workflow steps, making explicit which intermediate
step outputs will be returned from the workflow. All connections include
identifiers, which CWL document authors are encouraged to name
meaningfully, e.g., reference\_genome instead of input7.

CWL workflows form explicit data flows, as required for the particular
computational analysis. The connectivity between steps defines the
partial execution order. Parallel execution of steps is permitted and
encouraged whenever multiple steps have all of their inputs satisfied,
e.g., in Figure 1, find\_16S\_matches and find\_S5\_matches are at the
same data dependency level and can execute concurrently or sequentially
in any order. Additionally, a scatter construct allows the repeated
execution of a CWL step (perhaps overlapping in time, depending on the
resources available) where most of the inputs are the same except for
one or more inputs that vary. This is done without requiring the
modification of the underlying tool description. Starting with CWL
version 1.2, workflows can also conditionally skip execution of a (tool
or workflow) step, based upon a specified intermediate input or custom
boolean evaluation. Combining these features allows for a flexible
branch mechanism that allows workflow engines to calculate data
dependencies before the workflow starts, and thus retains the
predictability of the data flow paradigm.

In contrast to hard-coded approaches that rely on implicit filepaths
particular for each workflow, CWL workflows are more flexible, reusable,
and portable (which enables scalability). The use in the CWL standards
of explicit runtime environments, combined with explicit inputs/outputs
to form the data flow, enables step reordering and explicit handling of
iterations. The same features enable scalable remote execution and, more
generally, flexible use of runtime environments. Moreover, individual
tool definitions from multiple workflows can be reused in any new
workflow.

CWL workflow descriptions are also future-proof. Forward compatibility
of CWL documents is guaranteed, as each CWL document declares which
version of the standards it was written for and minor versions do not
alter the required features of the major version. A stand-alone upgrader
\(^{18}\) can automatically upgrade CWL documents from one version to
the next, and many CWL-aware platforms will internally update
user-submitted documents at runtime.

\section{3.1 Execution of workflows in CWL
format}\label{execution-of-workflows-in-cwl-format}

CWL is a set of standards, not a particular software product to install,
purchase, or rent. The CWL standards need to be implemented to be
useful; a list of some implementations of the CWL standards is in Table
1. Workflow/tool runners that claim compliance with the CWL standards
are allowed significant flexibility in how and where they execute a
user's CWL documents as long as they fulfill the requirements written in
those documents. For example, they are allowed (and encouraged) to
distribute execution of a workflow across all available computers that
can fulfill the resource requirements specified by the user. Aspects of
execution not defined by the CWL standards include (web) APIs for
workflow execution and real-time monitoring.

For example details about when a step should be considered ready for
execution are available in §4 of CWL Workflow Description standard
\(^{19}\) but once all the inputs are available the exact timing is up
to the workflow engine itself.

Step execution may result in a temporary or permanent failure, as
defined in §4 of CWL Workflow Description standard \(^{20}\) . It is up
to the workflow engine to control any automatic attempts to recover from
failures, e.g., to re-execute a Workflow step. Most workflow engines
that implement the CWL standards offer the feature of attempting a
number of re-executions, as set by the user, before reporting permanent
failure.

The CWL community has developed the following optimizations without
requiring that users re-write their workflows to benefit:

\begin{enumerate}
\def\labelenumi{(\arabic{enumi})}
\tightlist
\item
  Automatic streaming of data inputs and outputs instead of waiting for
  all the data to be downloaded or uploaded (where those data inputs or
  outputs are marked with ``streamable: true'')\\
\item
  Workflow step placement based upon data location {[}18{]}, resource
  needs, and/or cost of data transfer {[}19{]}\\
\item
  The re-use of the results from previously computed steps, even from a
  different workflow, as long as the inputs are identical. This can be
  controlled by the user via the ``WorkReuse'' directive \(^{21}\) .
\end{enumerate}

Real world usage at scale: routinely CWL users and vendors report that
they analyze 5000 whole genome sequences in a single workflow execution;
one customer of a commercial vendor reported a successful run of a
workflow that contained an 8,000-wide step; the entire workflow had
25,000 container executions. By design, the CWL standards do not impose
any technical limitations to the size of files processed or to the
number of tasks run in parallel. The major scalability bottlenecks are
hardware-related --- not having enough machines with enough memory,
compute or disk space to process more and more data at a larger scale.
As these boundaries move in the future with technological advances, the
CWL standards should be able to keep up and not be a cause of
limitations.

\section{3.2 When is CWL not useful?}\label{when-is-cwl-not-useful}

The CWL standards were designed for a particular style of command-line
tool based data analysis. Therefore, the following situations are out of
scope and not appropriate (or possible) to describe using CWL syntax:

\begin{enumerate}
\def\labelenumi{(\arabic{enumi})}
\tightlist
\item
  Safe interaction with stateful (web) services\\
\item
  Real-time communication between workflow steps\\
\item
  Interactions with command line tools beside 1) constructing the
  command line and making available file inputs (both user provided and
  synthesized from other inputs just prior to execution) and 2)
  consuming the output of the tool once its execution is finished, in
  the form of files created/changed,
\end{enumerate}

the POSIX standard output and error streams, and the POSIX exit code of
the tool

\begin{enumerate}
\def\labelenumi{(\arabic{enumi})}
\setcounter{enumi}{3}
\tightlist
\item
  Advanced control-flow techniques beyond conditional steps\\
\item
  Runtime workflow graph manipulations: dynamically adding or removing
  new steps during workflow execution, beyond any predefined conditional
  step execution tests that are in the original workflow description\\
\item
  Workflows that contain cycles: ``repeat this step or sub-workflow a
  specific number of times'' or ``repeat this step or sub-workflow until
  a condition is met.'' \(^{22}\)\\
\item
  Workflows that need particular steps run at or during a specific
  day/time-frame
\end{enumerate}

\section{4 OPEN-SOURCE, OPEN STANDARDS, OPEN
COMMUNITY}\label{open-source-open-standards-open-community}

Given the numerous and diverse set of potential users, implementers, and
other stakeholders, we posit that a project like CWL requires the
combined development of code, standards, and community. Indeed, these
requirements were part of the foundational design principles for CWL
(Section 4.1); in the long run, these have fostered free and open source
software (Sidebar B, in Section 4.2), and a vibrant and active ecosystem
(Section 4.3).

\section{4.1 The CWL Principles}\label{the-cwl-principles}

The CWL project is based on a set of five principles:

Principle 1: The core of the project is the community of people who care
about its goals.

Principle 2: To achieve the best possible results, there should be few,
if any, barriers to participation. Specifically, to attract people with
diverse experiences and perspectives, there must be no cost to
participate.

Principle 3: To enable the best outcomes, project outputs should be used
as people see fit. Thus, the standards themselves must be licensed for
reuse, with no acquisition price.

Principle 4: The project must not favor any one company or group over
another, but neither should it try to be all things to all people. The
community decides.

Principle 5: The concepts and ideas must be tested frequently: tested
and functional code is the beginning of evaluating a proposal, not the
end.

In time, the CWL project-members learned that this approach is a
superset of the OpenStand Principles \(^{23}\) , a joint ``Modern
Paradigm for Standards'' promoted by the IAB, IEEE, IETF, Internet
Society, and W3C. The CWL project additions to the OpenStand Principles
are: (1) to keep participation free of cost, and (2) the explicit choice
of the Apache 2.0 license for all its text, conformance tests, and
reference implementations.

Necessary and sufficient: All these principles have proven to be
essential for the CWL project. For example, the free cost and open
source license (Principles 2 and 3) has enabled many implementations of
the CWL standards, several of which re-use different parts of the
reference implementation of the CWL standards (reference runner). Being
community-first (Principle 1) has led to several projects from
participants that are outside the CWL standards themselves; the most
important contributions have made their way back into the project
(Principle 4).

Table 1: Selected F/OSS workflow runners and platforms that implement
the CWL standards.

Implementation

Platform support

cwltool

Linux, macOS, Windows (via WSL 2)local execution only

Arvados

in the cloud on AWS, Azure and GCP,on premise \& hybrid clusters using
Slurm

Toil

AWS, Azure, GCP, Grid Engine, HTCondor,LSF, Mesos, OpenStack, Slurm,
PBS/Torquealso local execution on Linux, macOS,MS Windows (via WSL 2)

CWL-Airflow

Local execution on Linux, OS Xor via dedicated Airflow enabled cluster.

REANA

Kubernetes, CERN OpenStack,OpenStack Magnum

As part of Principle 5, contributors to the CWL project have developed a
suite of conformance tests for each version of the CWL standards. These
publicly available tests were critical to the CWL project's success:
they helped assess the reference implementation of the CWL standards
themselves; they provided concrete examples to early adopters; and they
enabled the developers and users of production implementations of the
CWL standards to confirm their correctness.

\section{4.2 Sidebar B: The CWL project and Free/Open Source Software
(F/OSS)}\label{sidebar-b-the-cwl-project-and-freeopen-source-software-foss}

4.2.1 Free and Open Source implementations of the CWL standards. 24

By 2021, the CWL standards have gained much traction and are currently
widely supported in practice. In addition to the implementations in
Table 1, Galaxy \([2]^{25}\) and Pegasus \([10]^{26}\) have
in-development support for the CWL standards as well.

Wide adoption benefits from our principles: The CWL standards include
conformance tests, but the CWL community does not yet test or certify
implementations of the CWL standards, or specific technology stacks.
Instead, the authors and service provides of workflow runners and
workflow management systems self-certify support for the CWL standards,
based on a particular technology configuration they deploy and maintain.

4.2.2 F/OSS tools and libraries for working with CWL format documents.
\(^{27}\)

CWL plugins for text/code editors exist for Atom, vim, emacs, Visual
Studio Code, IntelliJ, gedit, and any text editor that support the
``language server protocol'' \(^{28}\) standard.

There are tools to generate CWL syntax from Python (via argparse/click
or via functions), ACD \(^{29}\) , CTD \(^{30}\) , and annotations in
IPython Jupyter Notebooks. Libraries to generate and/or read CWL
documents exist in many languages: Python, Java, R, Go, Scala,
Javascript, Typescript, and C++.

\section{4.3 The CWL Ecosystem}\label{the-cwl-ecosystem}

Beyond the ratified initial and updated CWL standards released over the
last six years, the CWL community has developed many tools, software
libraries, connected specifications, and has shared CWL descriptions for
popular tools. For example, there are software development kits for both
Python \(^{31}\) and Java \(^{32}\) that are generated automatically
from the CWL schema; this allows programmers to load, modify, and save
CWL documents using an object oriented model that has direct
correspondence to the CWL standards themselves. CWL SDKs for other
languages are possible by extending the code generation routines
\(^{33}\) . (See Sidebar B in Section 4.2.2 for practical details.)

The CWL standards support well the acute need to reuse (and,
correspondingly, to share) information on workflow execution, and on
authoring and provenance. The CWLProv \(^{34}\) prototype was created to
show how existing standards \([4, 21, 25]\) can be combined to represent
the provenance of a specific execution of a CWL workflow \([20]\) .
Although, to-date, CWLProv has only been implemented in the CWL
reference runner, interest is high in additional implementation and
further development.

\section{5 CONCLUSION}\label{conclusion}

The problem of standardizing computational reuse is only increasing in
prominence and impact. Addressing this problem, various domains in
science, engineering, and commerce have already started to migrate to
workflows, but efforts focusing on the portability and even definition
of workflows remain scattered. In this work we raise awareness to this
problem and propose a community-driven solution.

The Common Workflow Language (CWL) is a family of standards for the
description of command line tools and of workflows made from these
tools. It includes many features developed in collaboration with the
community: support for software containers, resource requirements,
workflow-level conditional branching, etc. Built on a foundation of five
guiding principles, the CWL project delivers open standards, open-source
code, and an open community.

For the past six years, the community around CWL has developed
organically. Organizations looking to write, use, or fund data analysis
workflows based upon command-line tools should adopt or even require the
CWL standards, because the CWL standards offer a common yet reduced set
of capabilities that are both used in practice and implemented in many
popular workflow systems. CWL is further valuable because it is
supported by a large-scale community, diverse fields have already
adopted it, and its adoption is rapidly growing. Specifically,

\begin{enumerate}
\def\labelenumi{(\arabic{enumi})}
\tightlist
\item
  By using a reduced set of capabilities, the CWL standards limit the
  complexity encountered by users when they start to use it, and by
  operators when they have to implement it. (Feedback from the community
  indicates these are appreciated.)\\
\item
  By using declarative syntax, CWL allows users to specify workflows
  even if they do not know exactly where the workflows would (later)
  run.\\
\item
  The CWL project is governed in the public interest and produces freely
  available open standards. The CWL project itself is not a specific
  workflow management system, workflow runner, or vendor. This allows
  potential users, operators, and vendors, to avoid lock-in and be more
  flexible in the future.\\
\item
  By offering standards, the CWL project distinguishes itself especially
  for the complex interactions that appear in scientific and engineering
  collaborations. These interactions include defining workflows from
  many different tools (or steps), sharing workflows, long-term
  archiving, fulfilling requirements of regulators (e.g., US FDA),
  making workflow executions auditable and reproducible. (This is
  particularly useful in cooperative environments, where groups that
  compete with each other need to collaborate, or in scientific papers
  where the paper results can be reused very efficiently if the analysis
  is described in a CWL workflow with publicly available software
  containers for all steps.)\\
\item
  The CWL standards are already implemented, adopted, and used; with
  many production-grade implementations available as open source and
  with zero-cost. Thus, the different communities of users of the CWL
  standards already offer numerous workflow and tool descriptions. (This
  is akin to how the Python ecosystem of shared libraries, code, and
  recipes is already helpful.)
\end{enumerate}

To conclude: this is a call for others to embrace workflow thinking and
join the CWL community in creating and sharing portable and complete
workflow descriptions. With the CWL standards, the methods are included
and ready to (re)use!

\section{ACKNOWLEDGMENTS}\label{acknowledgments}

The CWL project is immensely grateful to the following self-identified
CWL Community members and their contributions to the project: Miguel
d'Arcangues Boland (Software, Bug Reports, Maintenance), Alain Domissy
(Conceptualization, Answering Questions, Tools), Andrey Kislyuk
(Software, Bug Reports), Brandi Davis-Dusenbery (Conceptualization,
Funding acquisition, Investigation, Project Administration, Resources,
Supervision, Business Development, Event Organizing, Talks), Niels Drost
(Funding Acquisition, Blogposts, Event Organizing, Tutorials, Talks),
Robert Finn (Data acquisition, Funding acquisition, Investigation,
Resources, Supervision), Michael Franklin (Software, Bug Reports,
Documentation, Event Organizing, Maintenance, Tools, Answering
Questions, Talks), (), Manabu Ishii (Blogposts, Documentation, Examples,
Event Organizing,

Maintenance, Tools, Answering Questions, Translation, Tutorials, Talks),
Sinisa Ivkovic (Software, Validation, Bug Reports, Tools), Alexander
Kanitz (Software, Business Development, Tools, Talks), Sehrish Kanwal
(Conceptualization, Formal Analysis, Investigation, Software,
Validation, Bug Reports, Blogposts, Content, Event Organizing,
Maintenance, Answering Questions, Tools, Tutorials, Talks, User
Testing), Andrey Kartashov (Conceptualization, Software, Validation,
Examples, Tools, Answering Questions), Farah Khan (Conceptualization,
Formal Analysis, Funding Acquisition, Software), Michael Kotliar
(Software, Validation, Bug Reports, Blogposts, Examples, Maintenance,
Answering Questions, Reviewed Contributions, Tools, Talks, User
Testing), Folker Meyer (Tools), Rupert Nash (Software, Bug Reports,
Talks, Videos), Maya Nedeljkovich (Software, Validation, Visualization,
Writing -- review \& editing, Bug Reports, Tools, Talks), Tazro Ohta
(Formal Analysis, Funding Acquisition, Resources, Validation, Bug
Reports, Blogposts, Content, Documentation, Examples, Event Organizing,
Answering Questions, Tools, Translation, Tutorials, Talks, User
Testing), Pjotr Prins (Blogposts, Packaging, Bug Reports), Manvendra
Singh (Software, Blogposts, Packaging, Tools, Reviewed Contributions),
Andrey Tovchigrechko (Conceptualization, Software, Bug Reports), Alan
Williams (Investigation), Denis Yuen (Software, Bug Reports,
Documentation, Tools), Alexander (Sasha) Wait Zaranek
(Conceptualization, Funding Acquisition), Sarah Wait Zaranek
(Conceptualization, Funding Acquisition, Project Administration,
Resources, Software, Accessibility, Bug Reports, Business Development,
Content, Examples, Event Organizing, Answering Questions, Tutorials,
Talks).

The contributions to the CWL project by the authors of this paper are:
Michael R. Crusoe (Conceptualization, Funding Acquisition,
Investigation, Project Administration, Resources, Software, Supervision,
Validation, Writing -- original draft, Bug Reports, Business
Development, Content, Documentation, Examples, Event Organizing,
Maintenance, Packaging, Answering Questions, Reviewing Contributions,
Tutorials, Talks), Sanne Abeln (Writing -- original draft, and review \&
editing, Conceptualization, Supervision, Funding acquisition), Alexandru
Iosup (Writing -- original draft, and review, editing,
Conceptualization, Supervision, Funding acquisition), Peter Amstutz
(Conceptualization, Formal Analysis, Funding Acquisition, Methodology,
Project Administration, Resources, Software, Supervision, Validation,
Visualization, Bug Reports, Business Development, Content,
Documentation, Examples, Event Organizing, Maintenance, Packaging,
Answering Questions, Reviewed Contributions, Security, Tools, Tutorials,
Talks), Nebojša Tijanić (Conceptualization, Software), Hervé Ménager
(Conceptualization, Funding Acquisition, Investigation, Methodology,
Project Administration, Resources, Software, Supervision, Bug Reports,
Business Development, Examples, Event Organizing, Maintenance, Tools,
Tutorials, Talks), Stian Soiland-Reyes (Conceptualization, Funding
Acquisition, Investigation, Project Administration, Resources, Software,
Supervision, Validation, User Testing, Writing -- review and editing,
Bug Reports, Blogposts, Business Development, Content, Documentation,
Examples, Event Organizing, Packaging, Tools, Answering Questions,
Reviewed Contributions, Security, Tutorials, Talks, Videos), Bogdan
Gavrilovic (Conceptualization, Software, Validation, Bug Reports,
Blogposts, Maintenance, Tools, Answering Questions, Reviewed
Contributions, User Testing), Carole A. Goble (Conceptualization,
Funding Acquisition, Resources, Supervision, Audio, Business
Development, Content, Examples, Event Organizing, Tools, Talks, Videos).

Funding acknowledgements: European Commission grants BioExcel-2 (SSR)
H2020-INFRAEDI-02-2018 823830
(https://cordis.europa.eu/project/id/823830), BioExcel (SSR)
H2020-EINFRA-2015-1 675728 (https://cordis.
europa.eu/project/id/675728), EOSC-Life (SSR) H2020-INFRAEOSC-2018-2
824087 (https://cordis.europa.eu/project/id/824087), EOSCPilot (MRC)
H2020-INFRADEV-2016-2 739563
(https://cordis.europa.eu/project/id/739563), IBISBA 1.0 (SSR)
H2020-INFRAIA-2017-1-two-stage 730976
(https://cordis.europa.eu/project/id/730976), ELIXIR-EXCELERATE (SSR,
HM) H2020-INFRADEV-1-2015-1 676559
(https://cordis.europa.eu/project/id/676559), ASTERICS (MRC)
INFRADEV-4-2014-2015 653477
(https://cordis.europa.eu/project/id/653477). ELIXIR the research
infrastructure for life-science data, Interoperability Platform
Implementation Study (MRC). 2018-CWL
(https://elixir-europe.org/about-us/commissioned-services/cwl-2018).
Various universities have also co-sponsored this project; we thank Vrije
Universiteit of Amsterdam, the Netherlands, where the first three
authors have their primary affiliation.

\section{REFERENCES}\label{references}

{[}1{]} 2020. IEEE Standard for Bioinformatics Analyses Generated by
High-Throughput Sequencing (HTS) to Facilitate Communication. Technical
Report 2791-2020. 1--16 pages.
https://doi.org/10.1109/IEEESTD.2020.9094416 Conference Name: IEEE Std
2791-2020.\\
{[}2{]} Enis Afgan, Dannon Baker, Bérénice Batut, Marius van den Beek,
Dave Bouvier, Martin Čech, John Chilton, Dave Clements, Nate Coraor,
Björn A Grüning, Aysam Guerler, Jennifer Hillman-Jackson, Saskia
Hiltemann, Vahid Jalili, Helena Rasche, Nicola Soranzo, Jeremy Goecks,
James Taylor, Anton Nekrutenko, and Daniel Blankenberg. 2018. The Galaxy
platform for accessible, reproducible and collaborative biomedical
analyses: 2018 update. Nucleic Acids Research 46, W1 (July 2018),
W537--W544. https://doi.org/10.1093/nar/gky379\\
{[}3{]} Yadu Babuji, Anna Woodard, Zhuozhao Li, Daniel S. Katz, Ben
Clifford, Rohan Kumar, Lukasz Lacinski, Ryan Chard, Justin M. Wozniak,
Ian Foster, Michael Wilde, and Kyle Chard. 2019. Parsl: Pervasive
Parallel Programming in Python. In Proceedings of the 28th International
Symposium on High-Performance Parallel and Distributed Computing (HPDC
'19). Association for Computing Machinery, New York, NY, USA, 25--36.
https://doi.org/10.1145/3307681.3325400\\
{[}4{]} Khalid Belhajjame, Jun Zhao, Daniel Garijo, Matthew Gamble,
Kristina Hettne, Raul Palma, Eleni Mina, Oscar Corcho, José Manuel
Gómez-Pérez, Sean Bechhofer, Graham Klyne, and Carole Goble. 2015. Using
a suite of ontologies for preserving workflow-centric research objects.
Journal of Web Semantics 32 (May 2015), 16--42.
https://doi.org/10.1016/j.websem.2015.01.003\\
{[}5{]} Tim Bell, Luca Canali, Eric Grancher, Massimo Lamanna, Gavin
McCance, Pere Mato Vila, Danilo Piparo, Jakub Moscicki, Alberto Pace,
Ricardo Brito Da Rocha, Tibor Simko, Tim Smith, and Enric Tejedor
Saavedra. 2017. Web-based Analysis Services Report. Technical Report
CERN-IT-Note-2018-004. CERN, Geneva, Switzerland.
http://cds.cern.ch/record/2315331/\\
{[}6{]} Michael R. Berthold, Nicolas Cebron, Fabian Dill, Thomas R.
Gabriel, Tobias Kötter, Thorsten Meinl, Peter Ohl, Kilian Thiel, and
Bernd Wiswedel. 2009. KNIME - the Konstanz information miner: version
2.0 and beyond. ACM SIGKDD Explorations Newsletter 11, 1 (Nov.~2009),
26--31. https://doi.org/10.1145/1656274.1656280\\
{[}7{]} Peter Couvares, Tevfik Kosar, Alain Roy, Jeff Weber, and Kent
Wenger. 2007. Workflow Management in Condor. In Workflows for e-Science:
Scientific Workflows for Grids, Ian J. Taylor, Ewa Deelman, Dennis B.
Gannon, and Matthew Shields (Eds.). Springer, London, 357--375.
https://doi.org/10.1007/978-1-84628-757-2\_22\\
{[}8{]} Víctor Cuevas-Vicenttín, Saumen Dey, Sven Köhler, Sean Riddle,
and Bertram Ludäscher. 2012. Scientific Workflows and Provenance:
Introduction and Research Opportunities. Datenbank-Spektrum 12, 3
(Nov.~2012), 193--203. https://doi.org/10.1007/s13222-012-0100-z\\
{[}9{]} Luis de la Garza, Johannes Veit, Andras Szolek, Marc Röttig,
Stephan Aiche, Sandra Gesing, Knut Reinert, and Oliver Kohlbacher. 2016.
From the desktop to the grid: scalable bioinformatics via workflow
conversion. BMC Bioinformatics 17, 1 (March 2016), 127.
https://doi.org/10.1186/s12859-016-0978-9\\
{[}10{]} Ewa Deelman, Karan Vahi, Gideon Juve, Mats Rynge, Scott
Callaghan, Philip J. Maechling, Rajiv Mayani, Weiwei Chen, Rafael
Ferreira da Silva, Miron Livny, and Kent Wenger. 2015. Pegasus, a
workflow management system for science automation. Future Generation
Computer Systems 46 (May 2015), 17--35.
https://doi.org/10.1016/j.future.2014.10.008\\
{[}11{]} Dror G. Feitelson. 2015. From Repeatability to Reproducibility
and Corroboration. ACM SIGOPS Operating Systems Review 49, 1
(Jan.~2015), 3--11. https://doi.org/10.1145/2723872.2723875\\
{[}12{]} Peter Georgeson, Anna Syme, Clare Sloggett, Jessica Chung,
Harriet Dashnow, Michael Milton, Andrew Lonsdale, David Powell, Torsten
Seemann, and Bernard Pope. 2019. Bionitio: demonstrating and
facilitating best practices for bioinformatics command-line software.
*GigaScience* 8, giz109 (Sept.~2019).
https://doi.org/10.1093/gigascience/giz109

{[}13{]} Pedro Gonçalves. 2020. OGC Earth Observations Applications
Pilot: Terradue Engineering Report. OGC Public Engineering Report OGC
20-042. Open Geospatial Consortium.
http://docs.opengeospatial.org/per/20-042.html\\
{[}14{]} The Austin Group. 2008. POSIX.1-2008 (IEEE Std 1003.1 \(^{™}\)
-2008 and The Open Group Technical Standard Base Specifications, Issue
7). Technical Report. Austin Group.
https://pubs.opengroup.org/onlinepubs/9699919799.2008edition/\\
{[}15{]} Michael R. Gryk and Bertram Ludäscher. 2017. Workflows and
Provenance: Toward Information Science Solutions for the Natural
Sciences. Library trends 65, 4 (2017), 555--562.
https://doi.org/10.1353/lib.2017.0018\\
{[}16{]} Andrea Guarracino, Peter Amstutz, Thomas Liener, Michael
Crusoe, Adam Novak, Erik Garrison, Tazro Ohta, Bonface Munyoki, Danielle
Welter, Sarah Wait Zaranek, Alexander (Sasha) Wait Zaranek, and Pjotr
Prins. 2020. COVID-19 PubSeq: Public SARS-CoV-2 Sequence Resource.
https://bcc2020.sched.com/event/coLw/covid-19-pubseq-public-sars-cov-2-sequence-resource\\
{[}17{]} Peter Ivie and Douglas Thain. 2018. Reproducibility in
Scientific Computing. Comput. Surveys 51, 3 (July 2018), 63:1--63:36.
https://doi.org/10.1145/3186266\\
{[}18{]} Fan Jiang, Claris Castillo, and Stan Ahalt. 2019. TR-19-01: A
Cloud-Agnostic Framework for Geo-Distributed Data-Intensive
Applications. Technical Report TR-19-01. RENCI, University of North
Carolina at Chapel Hill. 10 pages.
https://renci.org/technical-reports/tr-19-01/\\
{[}19{]} Fan Jiang, Kyle Ferriter, and Claris Castillo. 2019. PIVOT:
Cost-Aware Scheduling of Data-Intensive Applications in a Cloud-Agnostic
System. Technical Report TR-19-02. RENCI, University of North Carolina
at Chapel Hill. 8 pages. https://renci.org/technical-reports/tr-19-02/\\
{[}20{]} Farah Zaib Khan, Stian Soiland-Reyes, Richard O. Sinnott,
Andrew Lonie, Carole Goble, and Michael R. Crusoe. 2019. Sharing
interoperable workflow provenance: A review of best practices and their
practical application in CWLProv. GigaScience 8, 11 (Nov.~2019).
https://doi.org/10.1093/gigascience/giz095\\
{[}21{]} J. Kunze, J. Littman, E. Madden, J. Scancella, and C. Adams.
2018. The BagIt File Packaging Format (V1.0). RFC 8493. RFC Editor.
https://www.rfc-editor.org/info/rfc8493 ISSN: 2070-1721.\\
{[}22{]} Tom Landry. 2020. OGC Earth Observation Applications Pilot:
CRIM Engineering Report. OGC Public Engineering Report OGC 20-045. Open
Geospatial Consortium. http://docs.opengeospatial.org/per/20-045.html\\
{[}23{]} Jessica W. Lau, Erik Lehnert, Anurag Sethi, Raunaq Malhotra,
Gaurav Kaushik, Zeynep Onder, Nick Groves-Kirkby, Aleksandar Mihajlovic,
Jack DiGiovanna, Mladen Srdic, Dragan Bajcic, Jelena Radenkovic,
Vladimir Mladenovic, Damir Krstanovic, Vladan Arsenijevic, Djordje
Klisic, Milan Mitrovic, Igor Bogicevic,

Deniz Kural, and Brandi Davis-Dusenbery. 2017. The Cancer Genomics
Cloud: Collaborative, Reproducible, and Democratized---A New Paradigm in
Large-Scale Computational Research. Cancer Research 77, 21 (Oct.~2017),
e3--e6. https://doi.org/10.1158/0008-5472.can-17-0387\\
{[}24{]} Jae-Hak Lee, Hana Yi, and Jongsik Chun. 2011. rRNASelector: A
computer program for selecting ribosomal RNA encoding sequences from
metagenomic and metatranscriptomic shotgun libraries. The Journal of
Microbiology 49, 4 (Sept.~2011), 689.
https://doi.org/10.1007/s12275-011-1213-z\\
{[}25{]} Paolo Missier, Khalid Belhajjame, and James Cheney. 2013. The
W3C PROV family of specifications for modelling provenance metadata. In
Proceedings of the 16th International Conference on Extending Database
Technology (EDBT '13). Association for Computing Machinery, New York,
NY, USA, 773--776. https://doi.org/10.1145/2452376.2452478\\
{[}26{]} Alex L Mitchell, Alexandre Almeida, Martin Beracochea, Miguel
Boland, Josephine Burgin, Guy Cochrane, Michael R Crusoe, Varsha Kale,
Simon C Potter, Lorna J Richardson, Ekaterina Sakharova, Maxim
Scheremetjew, Anton Korobeynikov, Alex Shlemov, Olga Kunyavskaya, Alla
Lapidus, and Robert D Finn. 2020. MGNify: the microbiome analysis
resource in 2020. Nucleic Acids Research 48, D1 (Jan.~2020), D570--D578.
https://doi.org/10.1093/nar/gkz1035\\
{[}27{]} H. Oliver, M. Shin, D. Matthews, O. Sanders, S. Bartholomew, A.
Clark, B. Fitzpatrick, R. v Haren, R. Hut, and N. Drost. 2019. Workflow
Automation for Cycling Systems: The Cyc Workflow Engine. Computing in
Science Engineering (2019), 1--1.
https://doi.org/10.1109/MCSE.2019.2906593 00000.\\
{[}28{]} Jeffrey M. Perkel. 2019. Workflow systems turn raw data into
scientific knowledge. Nature 573 (Sept.~2019), 149--150.
https://doi.org/10.1038/d41586-019-02619-z\\
{[}29{]} Torsten Seemann. 2013. Ten recommendations for creating usable
bioinformatics command line software. GigaScience 2, 2047-217X-2-15
(Dec.~2013). https://doi.org/10.1186/2047-217X-2-15\\
{[}30{]} Ingo Simonis. 2020. OGC Earth Observation Applications Pilot:
Summary Engineering Report. OGC Public Engineering Report OGC 20-073.
Open Geospatial Consortium. https://docs.ogc.org/per/20-073.html\\
{[}31{]} Ronald C. Taylor. 2010. An overview of the
Hadoop/MapReduce/HBase framework and its current applications in
bioinformatics. BMC Bioinformatics 11, 12 (Dec.~2010), S1.
https://doi.org/10.1186/1471-2105-11-S12-S1\\
{[}32{]} W. J. S. M. van Wezenbeek, H. J. J. Touwen, A. M. C. Versteeg,
and A. J. M. van Wesenbeeck. 2017. Nationaal plan open science.
Technical Report. Ministerie van Onderwijs, Cultuur en Wetenschap.
https://doi.org/10.4233/uuid:9e9fa82e-06c1-4d0d-9e20-5620259a6c65

\end{document}
